%----------------------------%
\section{Introduction}\label{sec:1}
%----------------------------%
%% Start to write the introduction
We consider linear systems of the following form:
%
\begin{equation}\label{sec:1:eq:FLS}
    AB\boldsymbol{x}= \boldsymbol{b},
\end{equation}
%
where $A\in \mathbb{R}^{N\times N}$ and $B\in \mathbb{R}^{N\times N}$ are assumed to be large and sparse nonsingular matrices, and $\boldsymbol{b} \in \mathbb{R}^N$ is a right-hand side vector. Linear system \eqref{sec:1:eq:FLS}  arises in the subproblems of product eigenvalue problems \cite{WA05} and a Pad\'e approximation to matrix function \cite[p.204]{SO23}. When $A\in \mathbb{R}^{m \times k}$, $B \in \mathbb{R}^{k \times n}$, the problem is a generalization of the normal
equation of least square problem $A^\top A\boldsymbol{x}=\boldsymbol{b}$, where $A\in \mathbb{R}^{m\times k}$, $\boldsymbol{b}\in \mathbb{R}^{k}$. This generalization of the normal equation is discussed in \cite{DU24}. However, this situation is different from our problem setting.
\par
To the best of our knowledge, there are few articles on numerical methods for linear system \eqref{sec:1:eq:FLS}. In our article, three numerical methods for linear systems \eqref{sec:1:eq:FLS} are considered.  Each method implements numerical methods for linear system $A\boldsymbol{x}=\boldsymbol{b}$ (e.g. direct methods, stationary iterative methods and Krylov subspace methods) by one implementing approach, leading to numerical methods for linear system \eqref{sec:1:eq:FLS}. We present the
three implementing approaches as follows.

\par
\begin{enumerate}
    \item The first approach considers the product of two factors $ A , B$ as one operator $ C $, thus the linear system $C\boldsymbol{x}=\boldsymbol{b}$ is solved.
    \item  The second approach first solves $A\boldsymbol{y}=\boldsymbol{b}$, then solves $B\boldsymbol{x}=\boldsymbol{y}$.
    \item  The third approach solves the augmented system
\begin{equation*}%\label{sec:2:eq:ALS}
    \mathcal{A}\Tilde{\boldsymbol{x}}=\Tilde{\boldsymbol{b}},\quad
\mathcal{A} = \begin{bmatrix}
    I & -B\\
    A & O
\end{bmatrix},
\Tilde{\boldsymbol{x}} =\begin{bmatrix}
    \boldsymbol{y}\\
    \boldsymbol{x}
\end{bmatrix},
\Tilde{\boldsymbol{b}}=\begin{bmatrix}
    \boldsymbol{0}\\
    \boldsymbol{b}
\end{bmatrix},
\end{equation*}
the solution for \eqref{sec:1:eq:FLS} is $\boldsymbol{x}$.
\end{enumerate}
\par
Our research is motivated by a question raised in \cite{DU24}. This article mentioned the second approach under situations where iterative methods are used. It is pointed out in \cite[p.1]{DU24} that this approach can be not accurate if the firstly solved system $A\boldsymbol{y}=\boldsymbol{b}$ terminates too early. The question naturally comes that, when should the firstly solved system $A\boldsymbol{y}=\boldsymbol{b}$ and the secondly solved system
$B\boldsymbol{x}=\boldsymbol{y}$ terminate, such that the accuracy of $AB\boldsymbol{x}=\boldsymbol{b}$ can be controlled? To answer this question, we propose a method to set the stopping criteria based on our theoretical analysis, and verify this method by numerical experiment, which is our main contribution.
Furthermore, we conduct experimental studies on the three implementing approaches with Krylov subspace methods (refer to \cite{SO23}, \cite{ME20} for details). This type of method is usually used for large and sparse matrices.
\par
The remainder of this paper is structured as follows: Section 2 first gives a detailed introduction on the second approach, then offeres a theoretical analysis on the accuracy of this approach. Then, based on the theoretical analysis, we develop our method for setting the stopping criteria, by which the accuracy of solution $\boldsymbol{x}$ is controlled. Section 3 conducts numerical experiments. The first experiment verifies the stopping criteria proposed in Section 2, while the
second experiment tests the three approaches on large and sparse matrices. Section 4 gives our conclusions.
\par
Throughout the paper $\|\cdot \|$ denotes  2-norm, i.e.,  $\|\cdot \|_2$. 

%----------------------------%
\section{Theoretical analysis on the accuracy of the second approach and a method to control the accuracy }\label{sec:2}
%----------------------------%
%----------------------------%
% \subsection{ORI approach (the first approach)}\label{subsec:ori}
%----------------------------%

%----------------------------%
% \subsection{SUB approach (the second approach)}\label{sec:2:sub:2}
%----------------------------%
\par
In this section, we introduce the seconde approach which features splitting linear system \eqref{sec:1:eq:FLS} into two subsystems. After that, we offer a theoretical analysis of the accuracy of this approach, and present our method to control the accuracy of this approach by setting the stopping criteria.
\par
The second approach exploits the following equavilent form of \eqref{sec:1:eq:FLS}:
 \begin{equation}\label{sec:2:eq:SLS}
    \begin{cases}
        A\boldsymbol{y}  =\boldsymbol{b},\\
        B\boldsymbol{x}  =\boldsymbol{y}.
    \end{cases}
\end{equation}
The solution $\boldsymbol{x}$ is obtained by sequentially solving the above two linear subsystems $A\boldsymbol{y}=\boldsymbol{b}$ and $B\boldsymbol{x}=\boldsymbol{y}$. Implementing methods should be combined with a linear solver, take, for example, Krylov subspace methods combined with the second approach. In the example method, Krylov subspace methods are applied first to the subsystem $A\boldsymbol{y}=\boldsymbol{b}$ with stopping criterion $ \alpha_A$, then to the subsystem
$B\boldsymbol{x}=\boldsymbol{y}$ with stopping criterion $ \alpha_B$. We evaluate the accuracy of the solution $\tilde{\boldsymbol{x}}$ for a linear system $A\boldsymbol{x}=\boldsymbol{b}$ by true relative residual $\varepsilon =
\frac{\left\|\boldsymbol{b}-A\tilde{\boldsymbol{x}}\right\|}{\left\|\boldsymbol{b}\right\|}$. The solution $\tilde{\boldsymbol{x}}$ could be less accurate if the true relative residual $ \varepsilon$ is larger. 
\par
From the definition above, we notice that $\alpha_A$ and  $\alpha_B$ are stopping criteria for the subsystems $A\boldsymbol{y}=\boldsymbol{b}$ and $B\boldsymbol{x}=\boldsymbol{y}$, instead of linear system \eqref{sec:1:eq:FLS}. It is worth considering the relation between the accuracies of the two subsystems \eqref{sec:2:eq:SLS} and the accuracy of linear system \eqref{sec:1:eq:FLS}.
The following theorem discusses the relation.
\par
\begin{thm}\label{thm:main}
    Let $\tilde{\boldsymbol{y}}$ be an approximation to the accurate solution $\boldsymbol{y}$ of
    $A\boldsymbol{y}=\boldsymbol{b}$, and $\tilde{\boldsymbol{x}} $ be an approximation to the accurate solution $\boldsymbol{x}$ of $B\boldsymbol{x}=\tilde{\boldsymbol{y}}$. Then the true relative residual $\varepsilon = \frac{\left\|\boldsymbol{b}-AB
            \tilde{\boldsymbol{x}}\right\|}{\left\|\boldsymbol{b}\right\|}
$ of solution $\tilde{\boldsymbol{x}}$ for \eqref{sec:1:eq:FLS} can be bounded by the following inequalities:
    %
    \begin{equation}\label{sec:2:eq:EA}
        \begin{aligned}
                    \frac{\left\|\boldsymbol{b}-AB \tilde{\boldsymbol{x}}\right\|}     
            {\left\|\boldsymbol{b}\right\|}
            &\leq 
            \kappa(A)\frac{\left\|\tilde{\boldsymbol{y}}-B
            \tilde{\boldsymbol{x}}\right\|}{\left\|\tilde{\boldsymbol{y}}\right\|}
            \frac{\left\|A\tilde{\boldsymbol{y}}\right\|}
            {\left\|\boldsymbol{b}\right\|}
            + \frac{\left\|\boldsymbol{b}-A
            \tilde{\boldsymbol{y}}\right\|}{\left\|\boldsymbol{b}\right\|}
            %\label{sec:2:eq:lls}
            .\\
        \end{aligned}
%                %.
    \end{equation}

\begin{proof}
    From triangular inequality $\left\|\boldsymbol{b}-AB\tilde{\boldsymbol{x}}\right\| \leq 
    \left\|\boldsymbol{b}-A\tilde{\boldsymbol{y}}\right\|+\left\|A(\tilde{\boldsymbol{y}}-B\tilde{\boldsymbol{x}})\right\|$, we easily prove the first inequality. For the second inequality, let $\sigma_{\min}$ and $\sigma_{\max}$ be the smallest and largest singular values of $A$, respectively. For any $\boldsymbol{y}\in \mathbb{R}^N$, we have:
    \[
    \sigma_{\min}\left\|\boldsymbol{y}\right\|\leq\left\|A\boldsymbol{y}\right\|
    \leq \sigma_{\max}\left\|\boldsymbol{y}\right\|
    ,\] 
    then $\frac{\left\|A(\tilde{\boldsymbol{y}}-
    B\tilde{\boldsymbol{x}})\right\|}{\left\|A\tilde{\boldsymbol{y}}\right\|}
    \leq \frac{\sigma_{\max}}{\sigma_{\min}}
    \frac{\left\|\tilde{\boldsymbol{y}}- B\tilde{\boldsymbol{x}}\right\|}
    {\left\|\tilde{\boldsymbol{y}}\right\|}.$ Notice $ \kappa(A) = 
    \frac{\sigma_{\max}}{\sigma_{\min}}
    $, the second inequality is thus proved.
\end{proof}
\end{thm}

The following two corollaries discuss special cases of inequality \eqref{sec:2:eq:EA} under given conditions, by which we develop the method for setting the stopping criteria $\alpha_A$, $\alpha_B$ to achieve controlled accuracy.

\par

\begin{col}\label{col:1}Let $ \varepsilon_A=\frac{\left\|\boldsymbol{b}-A
            \tilde{\boldsymbol{y}}\right\|}{\left\|\boldsymbol{b}\right\|}
$ be the true relative residual of solution $\tilde{\boldsymbol{y}}$ for $A\boldsymbol{y}=\boldsymbol{b}$. For any $\delta >0$, when $\varepsilon_A \leq \delta$, then $\left| \frac{\|A\tilde{\boldsymbol{y}}\|}{\|\boldsymbol{b}\|} -1 \right| \leq \delta$, thus $\frac{\left\|A\tilde{\boldsymbol{y}}\right\|}{\left\|\boldsymbol{b}\right\|} \leq 1+\delta$, the inequality \eqref{sec:2:eq:EA} becomes:
    \begin{equation}\label{sec2:eq:EB}
        \begin{aligned}
                    \frac{\left\|\boldsymbol{b}-AB \tilde{\boldsymbol{x}}\right\|}     
            {\left\|\boldsymbol{b}\right\|}& 
            &\leq
            \kappa(A)(1+\delta)\frac{\left\|\tilde{\boldsymbol{y}}-B
            \tilde{\boldsymbol{x}}\right\|}{\left\|\tilde{\boldsymbol{y}}\right\|}
            + \frac{\left\|\boldsymbol{b}-A
            \tilde{\boldsymbol{y}}\right\|}{\left\|\boldsymbol{b}\right\|}
            %\label{Sec2:eq:lls}
            .\\
        \end{aligned}
%                %.
    \end{equation}
    \begin{proof}
        From triangular inequality, $1-\frac{\left\|A\tilde{\boldsymbol{y}}\right\|}{\left\|\boldsymbol{b}\right\|} \leq \varepsilon_A $, $\frac{\left\|A\tilde{\boldsymbol{y}}\right\|}{\left\|\boldsymbol{b}\right\|}-1 \leq \varepsilon_A $, thus $\left| \frac{\|A\tilde{\boldsymbol{y}}\|}{\|\boldsymbol{b}\|} -1 \right| \leq \varepsilon_A \leq \delta$.
    \end{proof}
\end{col}
\par
In Corollary \ref{col:1}, we analyse the value of $\frac{\left\|A\tilde{\boldsymbol{y}}\right\|}{\left\|\boldsymbol{b}\right\|}$ and discover that as  $\varepsilon_A$ decreases, $\frac{\left\|A\tilde{\boldsymbol{y}}\right\|}{\left\|\boldsymbol{b}\right\|}$  can be arbitrarily close to 1. Based on the analysis, we can bound the true relative residual $\varepsilon
\leq \alpha$ by the true relative residuals of the subsystems in the following corollary.
\par
\begin{col}\label{col:2}
    Let $\varepsilon_B=\frac{\left\|\tilde{\boldsymbol{y}}-B
            \tilde{\boldsymbol{x}}\right\|}{\left\|\tilde{\boldsymbol{y}}\right\|}
$ be the true relative residual of solution $\tilde{\boldsymbol{x}}$ for $B\boldsymbol{x}=\tilde{\boldsymbol{y}}$. Assume $\delta > 0$ and  $\varepsilon_A \leq \delta$. For any $\alpha > 0$, if $ \varepsilon_A \leq \frac{\alpha}{2}, \varepsilon_B \leq \frac{ \alpha}{2(1+\delta)\kappa(A)} $, then $\varepsilon \leq \alpha$. In real computations, when $\tilde{\boldsymbol{y}}$ is a good approximation for $A\boldsymbol{y}=\boldsymbol{b}$, normally $\varepsilon_A\ll 1$, thus it can be assumed that
$\delta=1$. If $ \varepsilon_A \leq \frac{\alpha}{2}, \varepsilon_B \leq \frac{ \alpha}{4\kappa(A)} $, then $\varepsilon \leq \alpha$.

\begin{proof}
    From Corollary \ref{col:1}, $ \varepsilon \leq \kappa(A)(1+\delta) \varepsilon_B + \varepsilon_A$. It is assumed $ \varepsilon_A \leq \frac{\alpha}{2}, \varepsilon_B \leq \frac{ \alpha}{2(1+\delta)\kappa(A)} $, thus $\kappa(A)(1+\delta) \varepsilon_B + \varepsilon_A \leq \alpha$, then $ \varepsilon \leq \alpha$.
\end{proof}
\end{col}
    \par
    From the analysis in Corollary \ref{col:2}, it is reasonable if we set stopping criteria  $ \varepsilon_A \leq \frac{\alpha}{2}, \varepsilon_B \leq \frac{ \alpha}{4\kappa(A)} $ to control the accuracy, such that the true relative residual $\varepsilon \leq \alpha$, $\alpha$ is the given accuracy. As an application, we have developed our method to control the accuracy of the second approach by setting the stopping criteria $ \varepsilon_A \leq \frac{\alpha}{2}, \varepsilon_B \leq \frac{ \alpha}{4\kappa(A)} $. This method is summarized as the algotirhm \ref{alg:cg} below.

\begin{algorithm}
\caption{Solve linear system \eqref{sec:1:eq:FLS} by
    the separated approach}\label{alg:cg}
\begin{algorithmic}
\Require $A, B \in \mathbb{C}^{n \times n}, \bm{b} \in
      \mathbb{C}^n$, tolerance $\epsilon_{\text{tol}}$ for \eqref{sec:1:eq:FLS}
      \Ensure approximated solution $\bm{z}_m$ for linear
      system \eqref{sec:1:eq:FLS} s.t. $\| \bm{b}- \mathit{AB}
      \bm{z}_m \| \leq \epsilon_{\text{tol}} \| \bm{b} \|$.
      \State \textbf{Step 0}
      \State  $\epsilon_A = \frac{\epsilon_{\text{tol}}}{2}$, $\epsilon_B = \frac{\epsilon_{\text{tol}}}{(2 +
        \epsilon_{\text{tol}}) \kappa (A)}$
      \State \textbf{Step 1}
      \State Apply an iterative method to $A\bm{y}=\bm{b}$ to get $\bm{y}_m$ such that $\| \bm{b}- A\bm{y}_m
      \| \leq \epsilon_A \| \bm{b} \|$
      \State \textbf{Step 2}
      \State Apply an iterative method to $B\bm{z}=\bm{y}_m$ to get $\bm{z}_m$ such that $\| \bm{y}_m -
          B\bm{z}_m \| \leq \epsilon_B \| \bm{y}_m \|$
\end{algorithmic}
\end{algorithm}